% Generated draft: requires mathematical and source review.
\documentclass[11pt]{article}
\usepackage[margin=25mm]{geometry}
\usepackage{fontspec}
\setmainfont{Latin Modern Roman}
\usepackage{amsmath,amssymb,mathtools,unicode-math}
\setmathfont{Latin Modern Math}
\usepackage{enumitem,needspace,xurl,hyperref}
\setlength{\parindent}{0pt}
\setlength{\parskip}{5pt}
\setlength{\emergencystretch}{3em}
\newcommand{\N}{\mathbb{N}}
\newcommand{\Z}{\mathbb{Z}}
\newcommand{\Q}{\mathbb{Q}}
\newcommand{\R}{\mathbb{R}}
\newcommand{\C}{\mathbb{C}}
\newcommand{\F}{\mathbb{F}}
\newcommand{\A}{\mathbb{A}}
\newcommand{\PP}{\mathbb{P}}
\newcommand{\E}{\mathbb{E}}
\newcommand{\eps}{\varepsilon}
\newcommand{\dd}{\,\mathrm{d}}
\begin{document}
\section*{UnsolvedMath ID: 30000635}
\textbf{Source problem code:} catalogue entry
\subsection{Definitions and mathematical statement}
A convex body in $\R^3$ is a compact convex set with nonempty interior. For a unit direction $u\in S^2$, its parallel X-ray is the line-length function
\[
 X_uK(x)=\int_{\R}\mathbf1_K(x+tu)\,\dd t,
 \qquad x\in u^\perp.
\]
Thus the data give the lengths of all intersections with lines parallel to $u$, not merely the orthogonal projection or its area. The direction is unoriented: $u$ and $-u$ give the same data.

The question is whether every set of seven directions in the source's general position uniquely determines every convex body. With exact registered X-ray data, the requested implication is
\[
 \forall U=\{u_1,\ldots,u_7\}\text{ in general position}\
 \forall K,L\text{ convex bodies in }\R^3,
\]
\[
 \left[\forall i\in\{1,\ldots,7\}\ \forall x\in u_i^\perp:
                 X_{u_i}K(x)=X_{u_i}L(x)\right]
       \quad\Longrightarrow\quad K=L.
\]
The equality is within the class of convex bodies and in a common coordinate system; no unknown translations or rotations of separate data sets are introduced.

The source does not define general position. A common directional interpretation is that every three chosen direction vectors are linearly independent, or equivalently no three directions lie in one plane through the origin. Any additional geometric or arithmetic exclusions must be recovered rather than silently assumed. The source also does not state whether equality of X-rays is pointwise or almost everywhere; the displayed implication spells out exact pointwise data and flags that convention for review.

The quantifier is every admissible set of seven directions, not existence of one specially chosen set or an almost-everywhere choice. No smoothness, strict convexity, symmetry, or polytopal restriction is imposed on the bodies.
\subsection{Short English statement}
The question asks whether every general-position set of seven directions in three dimensions determines a convex body from all its line-intersection lengths. The universal direction quantifier is retained, while general position and data-equality conventions need confirmation.
\subsection{Sources}
\begin{itemize}
\item Local UnsolvedMath registry, permanent ID 30000635, code catalogue entry, ``Convex-Body Reconstruction from Seven Parallel X-Rays.'' Proposed year 2006; proposer unspecified.
\item \emph{Konvexgeometrie} (2007), Oberwolfach Reports, pp.\ 3321--3396, DOI 10.4171/owr/2006/56; \url{https://doi.org/10.4171/owr/2006/56}. The displayed report year 2007 differs from the DOI and proposed year 2006. This is the sole supplied literature reference.
\item Statement status corrected\_verified; the supplied note says only ambient-space notation and compound-title typography were normalized, retaining all quantifiers.
\item Literature triage checked 2026-08-21: open. It generically calls the item an unresolved problem/conjecture/direction and says the supplied check found no authoritative complete resolution. Exact general-position assumptions and current status require independent verification. Created 2026-08-25; updated 2026-09-16.
\item Registry curation by Ulam AI / UnsolvedMath Contributors, CC BY 4.0; the underlying report retains its supplied provenance. Unreviewed statement-only draft.
\end{itemize}
\end{document}
